\appendix

\section{The claim ledger}
\label{app:ledger}

Every numbered claim this project made, with its verification status as recorded at the time
of writing. Status is copied from the project's ledger rather than re-judged here. Claims
marked ``(1 seed)'' carry that qualifier in every sentence that makes them. Retracted
claims are in Appendix~\ref{app:retractions} and are not asserted anywhere in the body.

\begin{center}
\footnotesize
\begin{tabular}{l p{0.27\textwidth} p{0.34\textwidth} p{0.16\textwidth}}
  \toprule
  & claim & evidence & status \\
  \midrule
  C1 & algebra table is correct & \texttt{algebra.py::\_selfcheck}, primitive roots by explicit search & verified, self-checking \\
  C2 & non-regular $\mathcal{J}$-class counts $0,0,1,2,8$ & \texttt{j\_structure}; reproduces Prop.\ D.13 on their moduli & verified \\
  C3 & every modulus groks & $23$ moduli, up to $5$ seeds & verified \\
  C4 & $n=113$ replicates prior spectral statistics & $5$ seeds; key-frequency-count clause \textbf{withdrawn} & verified, 5 seeds \\
  C5 & analysis pipeline is correct & run on a third-party published checkpoint & verified externally \\
  C6 & clock survives prime powers & five cyclic prime powers & verified \\
  C7 & non-regular classes carry character structure & superseded by C33 & verified, \textbf{1 seed} \\
  C8 & key additive frequencies are CRT duals & $12$ moduli, $3$--$5$ seeds & verified \\
  C9 & $\mathrm{Gini}_{\text{add}}$ tracks zero-divisor density & $18$ moduli & verified \\
  C10 & from-scratch engine is correct & forward $4.4$e--$16$, gradients $2.68$e--$15$ & verified \\
  C11 & from-scratch engine groks & acceptance test & verified \\
  C12 & softmax collapse occurs on our engine & stamped artifact, $6.2$ orders & verified \\
  C13 & StableMax mitigates C12 & mechanism, not the earlier $10\times$ figure & verified (mechanism) \\
  C14 & critical dataset size is real here & $3$ seeds $\times$ $3$ moduli & verified \\
  C15 & engine reproduces prior addition result & all $8$ criteria & pass \\
  C17 & a frequency can be present yet vestigial & & verified, \textbf{1 seed} \\
  C18 & the addition circuit is a clean Clock & & verified, \textbf{1 seed} \\
  C21 & the circuit uses the embedding's key frequencies & $5$ seeds & verified \\
  C22 & clock frequencies are causally responsible at prime powers & $p<0.01$ in $5/5$ & verified, 5 seeds \\
  C23 & causal where there is no discrete log & $28/29$ analysable runs & verified \\
  C24 & $40$k steps is not converged at four moduli & PR retired at composite moduli & verified \\
  C25 & readout is invariant to the primitive root & exact to $10^{-10}$ & verified \\
  C26 & single-checkpoint Gini is noisy, sd $0.02$--$0.044$ & & verified \\
  C28 & the early CRT signal is CRT-specific & only magnitude discriminates & verified (specificity) \\
  C29 & all four moduli grok on the engine & $3$ seeds, float64 & verified \\
  C31 & the clock is in the neurons & $5$ seeds & verified \\
  C32 & causality survives all three initialisations & $24/24$ & verified \\
  C33 & every stratum runs a clock on its own local group & $32$ strata, $157$ measurements & verified, 5 seeds \\
  C34 & precision changes measured sparsity \emph{in PyTorch} & the $2\times2$ & verified \\
  C35 & $n=113$ is not special & three primes & verified \\
  C36 & the clock is causal at $2^{7}$ & $\hat{p}=10^{-4}$ in $2/2$; $2.77$--$476\times$ & verified \\
  C37 & $\omega(n)$ separates additive sparsity at matched zdd & $23$ moduli & \textbf{confirmatory} \\
  C38 & the same clock across strata, not one per stratum & $158/160$ pairs, null $10^{-4}$ & verified, 5 seeds$^{\dagger}$ \\
  C39 & key characters are what the network computes with & $3/3$ seeds, $n=113$, weight-space & verified \\
  C40 & the same at a non-square-free modulus, non-unit rows intact & $3/3$ seeds, $n=121$, weight-space & verified \\
  \bottomrule
\end{tabular}
\end{center}

\noindent
$^{\dagger}$C38 ships without a working failed-run control: the only two scorable
failed-run pairs in the project come from one model at test accuracy $0.7502$, and its own null
reads $p = 1.0$. Section~\ref{sec:strata-sameness} states this where the claim is made.

\paragraph{On the word \emph{causal} in this table.} These rows are the project's claim ledger,
quoted in the wording each claim was registered under. Throughout, \emph{causal} and
\emph{causally responsible} mean the logit-spectrum ablation of
Section~\ref{sec:methods-causal} (restricted and excluded losses against a permutation null),
and never weight ablation or activation patching, neither of which this paper performs.

\section{Retractions}
\label{app:retractions}

Five claims were made and withdrawn. Each is recorded with what killed it and the control that
would have caught it earlier. No paper in the corpus this one joins does this; we do it because
it is what makes the surviving claims checkable.

Three of the five died to the same shape, and it is worth stating once in symbols. A grouping
variable $X$ is compared against a response $Y$, but $X$ is structurally coupled to a size $S$
that $Y$ also tracks:
\begin{equation}
  \label{eq:confound}
  X \longrightarrow S \longrightarrow Y,
  \qquad\text{so } \rho(X, Y) \text{ is large while } \rho(X, Y \mid S) \text{ is not,}
\end{equation}
with $\rho(\cdot \mid \cdot)$ the partial correlation of \eqref{eq:partial}. The cheap general
defence, which caught all three, is to run the statistic on a control that has the size
structure and not the effect: shuffled labels, flat variance, a permuted grouping. If the
control reproduces the result, there is no result.

\paragraph{C19: $\varphi(n)$-normalised grokking time separates cyclic from non-cyclic
moduli.} \emph{Killed by:} a size confound in the normalisation itself. Raw grokking time
\emph{decreases} with $\varphi$ ($\rho = -0.665$), so dividing by $\varphi$ over-corrects: with
$S = \varphi(n)$ in \eqref{eq:confound}, $\rho(-\varphi, \text{steps}/\varphi) = +0.870$ over
$15$ moduli, an ordering manufactured entirely by the normalisation. The
cyclic moduli in the original sample were simply the three largest $\varphi$. \emph{The control
that would have caught it:} correlating the normalised statistic with the quantity it was
normalised by, before interpreting it.

\paragraph{C20: activation variance is graded by $\mathcal{J}$-class depth.} \emph{Killed
by:} a size confound between depth and block cell count, coupled through
$|J_{d}| = \varphi(n/d)$ in \eqref{eq:jclass}. $\rho(\text{cells},
\text{variance})$ runs $+0.856$ to $+0.949$, and a flat-variance \emph{noise} control scores
$\rho = +0.942$. \emph{The control:} measuring every block at an equal subsampled cell count,
and dropping blocks that cannot supply it rather than reporting them as zero.

\paragraph{C7b: regular $\mathcal{J}$-classes carry more structure than non-regular ones.}
\emph{Killed by:} a size confound between regularity and class size: $J_{1}$ is always
regular and, by \eqref{eq:jclass}, always the largest class. Size-matched within one model at $n = 120$, the only place such
pairs exist, the two overlap. \emph{The control:} the size-matched comparison, which is what
retired it.

\paragraph{C27: a network can lose a learned circuit.} \emph{Killed by:} right-censoring. The
claim was one run read at its final logged sample. Across $11$ long runs there are $19$
post-grok excursions below accuracy $0.90$, every one exactly one logging interval wide, $18$
of which recover. The nineteenth merely begins at the last step at which an observation
exists, and that run had recovered from a deeper excursion $4{,}600$ steps earlier. \emph{The control:}
classifying a run by a window rather than by its last row, which is now done everywhere.

\paragraph{C30: training precision changes measured sparsity.} \emph{Killed by:} a controlled
A/B that varied dtype inside our engine alone and returned the wrong sign, effect size
$-0.07$. It is superseded by C34, which is the same phenomenon as an \emph{interaction}:
precision is decisive in PyTorch and inert in our engine. C30's main-effect sentence is never
revived. \emph{The control:} running the variable. The claim had a $7/7$ correlation and a
mechanism that tied up three loose ends, which is the most persuasive shape a wrong claim can
take, and no amount of care in the statistics would have rescued it.

\section{Pre-registrations}
\label{app:prereg}

Every experiment whose criteria could be fixed in advance has a pre-registration committed
before the run, and the git timestamp is the evidence the criteria predated the result. Each is
named in the body beside the result it governs, with its commit. Two exceptions are stated
where they occur: the first gate's criteria, two of which were edited after failing
(Section~\ref{sec:calib}), and the $\omega(n)$ analysis, which is confirmatory because the
bands were seen before the pre-registration was written.

\section{Per-modulus and per-stratum tables}
\label{app:tables}

The full per-stratum table underlying Section~\ref{sec:strata}, and per-modulus spectral
statistics for all $23$ moduli, regenerate from the repository.

\section{Instruments}
\label{app:instruments}

The from-scratch autograd engine, its finite-difference test suite, the memory guards, and the
self-check suite that gates every claim-carrying script. One item is worth recording as a
limitation of its own: the engine's reported speedup is verified only for numerics, to
$2.68 \times 10^{-15}$; the wall-clock figure is unstamped and machine-dependent.
